\documentclass[12pt,a4paper]{article}
\usepackage{xurl}
\usepackage[slovak]{babel}
\usepackage[utf8]{inputenc}
\usepackage[T1]{fontenc}
\usepackage{times}
\usepackage[left=2.5cm,right=2cm,top=2.5cm,bottom=2.5cm]{geometry}
\usepackage{setspace}
\usepackage[backend=biber, style=iso-numeric, sorting=nty]{biblatex}
\usepackage{comment}
\usepackage{graphicx}
\usepackage{acro}
\begin{comment}
\DeclareAcronym{}{
    short={},
    long={}
}
\end{comment}

\DeclareAcronym{sav}{
    short={SAV},
    long={Slovenská Akadémia Vied}
}

\DeclareAcronym{scpi}{
  short = {SCPI},
  long  = {Standard Commands for Programmable Instruments}
}

\DeclareAcronym{rs-232}{
    short = {RS-232},
    long = {Recommended Standard 232}
}

\DeclareAcronym{gpib}{
  short = {GPIB},
  long  = {General Purpose Interface Bus}
}

\DeclareAcronym{ldls}{
  short = {LDLS},
  long  = {Laser-Driven Light Source®}
}

\DeclareAcronym{uv-c}{
    short= {UV-C},
    long= {Germicídne UV}
}

\DeclareAcronym{bnc}{
    short={BNC},
    long={Bayonet Neill–Concelman}
}

\DeclareAcronym{usb}{
    short={USB},
    long={Universal Serial Bus}
}

\DeclareAcronym{dll}{
    short={DLL},
    long={Dynamic-link library}
}

\DeclareAcronym{xml-rpc}{
    short={XML-RPC},
    long={Extensible Markup Language Remote Procedure Call}
}

\DeclareAcronym{mds}{
    short={MDS},
    long={Multi-Device Synchronization}
}

\addbibresource{literatura.bib}

\onehalfspacing

\usepackage{csquotes}
\begin{document}

\begin{titlepage}
    % Názov školy alebo inštitúcie (veľkosť 14, tučné)
    \begin{center}
        {\fontsize{14pt}{18pt}\selectfont \textbf{Elektrotechnický ústav SAV, v.v.i., Dúbravská cesta 9, 841 04 Bratislava}}
        
        \vspace{10cm}
        
        % Názov práce v strede strany, tučným písmom (veľkosť 16)
        {\fontsize{16pt}{20pt}\selectfont \textbf{Automatizácia optického meracieho pracoviska v jazyku Python}}
        
        \vspace{1cm}
        % Označenie práce (veľkosť 12)
        {\fontsize{12pt}{16pt}\selectfont Stredoškolská odborná činnosť}
    \end{center}

    \vspace*{\fill}

    % Miesto vľavo dole, meno a trieda vpravo dole (veľkosť 12)
    \noindent
    \begin{minipage}[t]{0.45\textwidth}
        \fontsize{12pt}{16pt}\selectfont
        Bratislava \\
        2026/2027
    \end{minipage}%
    \begin{minipage}[t]{0.55\textwidth}
        \raggedleft
        \fontsize{12pt}{16pt}\selectfont
        Maximilián Cenkner\\
        III.D
    \end{minipage}

\end{titlepage}

% -- POĎAKOVANIE --
\newpage
\vspace*{\fill}
\noindent\textbf{Poďakovanie}

\noindent Rád by som na tomto mieste poďakoval vedúcemu laboratória a tohto meracieho projektu na Elektrotechnickom ústave SAV, v.v.i., Ing. Mariánovi Vargovi, PhD. za realizáciu optickej stránky pracoviska a za fyzikálne a iné konzultácie spojené s pracoviskom a vypracovávaním tejto odbornej práce. Zároveň by som rád poďakoval Ing. Michalovi Sobotovi, PhD. za odborné konzultácie spojené s meraniami.
\vspace{2cm}
\thispagestyle{empty}


% -- OBSAH --
\newpage
\tableofcontents
\thispagestyle{empty}
\newpage

% -- ZOZNAM SKRATIEK A ZNAČIEK --
\section*{Zoznam skratiek a značiek}
\addcontentsline{toc}{section}{Zoznam skratiek a značiek}
\printacronyms[%
  heading=none, % Vypne podnadpisy, ak ich nechceš
  sort=true     % Automaticky zoradi skratky abecedne
]
\newpage

\section{Abstrakt}

% Kľúčové slová patria tesne pod abstrakt
\vspace{1cm}
\noindent\textbf{Kľúčové slová:} % Tu neskôr doplníš 3-5 kľúčových slov, napr. automatizácia, Python, optické meranie...
\newpage
\section{Úvod}

% MERACIE PRÍSTROJE
\begin{comment}
Identifikácia a primárna funkcia

Kľúčové technické parametre: Rozsahy, citlivosť, presnosť alebo iné fyzikálne limity. Uvádzaj len tie, ktoré sú relevantné pre tvoje meranie (napr. rozsah od 20 mA po 2 nA). Nemusíš prepisovať celý datasheet.
Komunikačné rozhranie: Keďže je tvoja SOČ zameraná na automatizáciu v Pythone, pri každom prístroji povinne uveď, cez aký port (USB, GPIB, RS-232, Ethernet) a akým jazykom/protokolom (SCPI príkazy, špecifické API výrobcu) s ním počítač komunikuje.

Úloha na pracovisku: Jasné prepojenie s experimentom. Prečo tam ten prístroj je a čo konkrétne robí?
\end{comment}

\newpage
\section{Použité meracie zariadenia}
\subsection{KEITHLEY 6487 Picoammeter/Voltage Source}
KEITHLEY 6487 Picoammeter/Voltage Source je pikoampérmeter s vysokým rozlíšením,
programovateľný s využitím \ac{scpi} príkazov cez port \ac{gpib} alebo \ac{rs-232}. Uvedený model 6487 dokáže merať prúd v 8 rôznych meracích rozsahoch, a to od 20~mA až po 2~nA. Pri meraniach sme použili tento prístroj aj ako zdroj $\pm 500$~V DC\footfullcite[s. 10]{keithley6487}.

Na pracovisku je prístroj pripojený k počítaču cez rozhranie \ac{gpib} pomocou prevodníka \ac{gpib}--\acs{usb} od firmy Agilent (Keysight), ktorý vyžaduje inštaláciu ovládačov Keysight IO Libraries Suite. V jazyku Python s prístrojom komunikujeme pomocou knižnice PyVISA\footfullcite{pyvisa}, cez ktorú odosielame \ac{scpi} príkazy a čítame namerané hodnoty. Pri meraní voltampérovej charakteristiky prístroj nastavuje napätie na vzorke a zároveň meria prúd, pri meraní responzivity slúži ako zdroj konštantného predpätia vzorky.

\subsection{Energetiq Fiber Coupled Laser-Driven Light Source EQ-99X-FC-S}
EQ-99X-FC-S je širokopásmový zdroj svetla s výstupnou optickou cestou a technológiou \ac{ldls}, ktorá pomocou zaostrených laserových lúčov vytvára plazmu, ktorá je zdrojom svetla. Svetelná šírka pásma vychádzajúca zo zdroja je v rozmedzí od 190~nm do 3500~nm, čo vyhovuje nášmu cieľu merať v rozmedzí \ac{uv-c} na svetelnom spektre. S prístrojom nie je možné komunikovať cez rozhranie PC\footfullcite{ldls}.

\newpage
\subsection{Monochromator LOT Quantum Design MSH-300M}
MSH-300 je monochromátor typu Czerny-Turner s ohniskovou vzdialenosťou 300~mm a 3 difrakčnými mriežkami slúžiacimi pre vyfiltrovanie zadanej vlnovej dĺžky. Monochromátor má 2 štrbiny slúžiace ako výstup a vstup pre svetlo, ktoré sú nastaviteľné s minimálnou a maximálnou šírkou od 10~\textmu m až do 10~mm. S prístrojom je možné komunikovať cez rozhranie \ac{usb}, a to pomocou \ac{dll} knižnice poskytnutej výrobcom\footfullcite{msh300m}.

Náš monochromátor je osadený mriežkami s hustotou 1800, 1200 a 600 vrypov/mm, optimalizovanými na vlnové dĺžky 250~nm, 500~nm a 1250~nm. Medzi mriežkami prepína automaticky podľa nastavenej vlnovej dĺžky, a to pri 400~nm a 1000~nm. Keďže knižnica \ac{dll} od výrobcu je 32-bitová, ovládanie monochromátora beží v samostatnom 32-bitovom procese jazyka Python, s ktorým hlavný program komunikuje pomocou protokolu \ac{xml-rpc}. Na pracovisku monochromátor vyberá zo širokopásmového svetla zdroja \ac{ldls} úzke spektrálne pásmo so zadanou vlnovou dĺžkou, ktorým sa osvetľuje meraná vzorka.

\subsection{Referenčné fotodetektory}
Pri meraní responzivity potrebujeme poznať optický výkon svetla dopadajúceho na vzorku. Na tento účel používame referenčné fotodetektory so známou (kalibrovanou) spektrálnou responzivitou. Z nameraného fotoprúdu referenčného detektora vypočítame optický výkon svetla, ktorý následne porovnáme s fotoprúdom meranej vzorky. Fotoprúd referenčného detektora meriame lock-in zosilňovačom MFLI.

\subsubsection{Thorlabs SM05PD2A Mounted Silicon Photodiode}
SM05PD2A je kremíková fotodióda s osadeným senzorom typu Thorlabs FDS010, schopná snímať a konvertovať svetelný signál na fotoprúd $I_{ph}$ v rozmedzí od 200~nm do 1100~nm. Diódu je možné zapojiť do meracieho zariadenia pomocou koaxiálneho kábla s koncovkou \ac{bnc}\footfullcite{thorlabs_sm05pd2a}.

\newpage
\subsubsection{Newport 918D-UV-OD3R Photodiode Sensor}
918D-UV-OD3R je kalibrovaný kremíkový fotodetektor so zvýšenou citlivosťou v UV oblasti, určený na meranie nízkych optických výkonov v rozsahu vlnových dĺžok od 200~nm do 1100~nm. Senzor má priemer 11,3~mm a je vybavený prepínateľným atenuátorom s optickou hustotou OD3\footfullcite{newport_918d}. K detektoru je dodaný kalibračný certifikát s krivkou responzivity (v A/W) s krokom 10~nm, a to pre meranie s atenuátorom aj bez neho. Kalibračnú krivku sme previedli do digitálnej podoby a v programe ju používame na prepočet nameraného fotoprúdu na optický výkon.

\subsection{Thorlabs ELL22 Motorized Neutral Density Filter Wheel Mount}
ELL22 je motorizovaný držiak filtra\footfullcite{thorlabs_ell22}. V rámci nášho merania je držiak osadený kontinuálnym filtrom NDC-50C-4\footfullcite{thorlabs_ndc_50c_4}, ktorý dokáže v reálnom čase upravovať výkon (intenzitu) prechádzajúceho svetla. V našom prípade so zariadením komunikujeme pomocou rozšírenia Thorlabs ELLC2\footfullcite{thorlabs_ellc2}, ktoré obsahuje rozhranie \ac{usb}, cez ktoré komunikácia prebieha. Na automatizáciu v Pythone sme využili knižnicu pylablib\footfullcite{pylablib}, ktorá obsahuje vysokoúrovňové funkcie, ktoré zjednodušujú automatizáciu.

\newpage
\subsection{Thorlabs FW102C 6-position Filter Wheel}
FW102C je motorizované filtrové koleso so 6 pozíciami pre optické filtre s priemerom 25~mm\footfullcite{thorlabs_fw102c}. Pre účely merania sú na jednotlivých pozíciách osadené nasledujúce dlhovlnné priepustné filtre (angl. \textit{longpass}):

\begin{itemize}
    \item Pozícia 1: Bez filtra, 100~\% prechod svetla
    \item Pozícia 2: FGL280 \O25~mm N-WG280 Colored Glass Filter, 280~nm Longpass
    \item Pozícia 3: FGL435 \O25~mm GG435 Colored Glass Filter, 435~nm Longpass
    \item Pozícia 4: FGL530 \O25~mm OG530 Colored Glass Filter, 530~nm Longpass
    \item Pozícia 5: FGL695 \O25~mm RG695 Colored Glass Filter, 695~nm Longpass
    \item Pozícia 6: SM1PL Externally SM1-Threaded Plug (100~\% zakrytie svetla)
\end{itemize}

Filtre slúžia na odfiltrovanie nežiaducich kratších vlnových dĺžok (napr. vyšších difrakčných rádov monochromátora). Program pri každej zmene vlnovej dĺžky automaticky natočí koleso na vhodnú pozíciu: pod 290~nm na pozíciu 1, od 290~nm na pozíciu 2, od 445~nm na pozíciu 3, od 545~nm na pozíciu 4 a od 720~nm na pozíciu 5. So zariadením komunikujeme cez rozhranie \ac{usb} (virtuálny sériový port) pomocou knižnice pylablib\footfullcite{pylablib}.

\newpage
\subsection{Thorlabs MC2000B Optical Chopper}
MC2000B je precízny optický svetelný chopper určený na moduláciu svetelného zväzku na danú frekvenciu. Kontinuálny zdroj svetla sa tak po prejdení chopperom mení na svetelné pulzy s určenou frekvenciou. Na MC2000B je možné osadiť rôzne typy clôn (angl. \textit{blade}) s danými parametrami. Pre účely merania sme použili clonu typu MC1F10HP s frekvenčným rozsahom $20$~Hz -- $1$~kHz\footfullcite{thorlabs_mc2000b}. So zariadením je možné komunikovať cez rozhranie \ac{usb} pomocou knižnice mc2000b v jazyku Python\footfullcite{mc2000b_docs}. Modulácia svetla umožňuje následné synchrónne meranie fotoprúdu lock-in zosilňovačmi MFLI na frekvencii choppera.

\subsection{Zurich Instruments MFLI Lock-In Amplifier 500~kHz}
MFLI je digitálny lock-in zosilňovač s frekvenčným rozsahom od DC do 500~kHz (po rozšírení až do 5~MHz). Na prednom paneli obsahuje jeden prúdový a jeden diferenciálny napäťový signálový vstup a vstupný signál digitalizuje 16-bitovým prevodníkom so vzorkovacou frekvenciou 60~MSa/s\footfullcite{zurich_mfli}. Lock-in zosilňovač z meraného signálu vyberá iba zložku s frekvenciou referenčného signálu (v našom prípade frekvenciou choppera), čím potláča šum a jednosmerné pozadie signálu, napríklad temný prúd fotodetektora.

Na pracovisku používame dva prístroje MFLI: jeden meria fotoprúd referenčného detektora a druhý fotoprúd meranej vzorky. Aby boli obidve merania časovo zosúladené, prístroje sú synchronizované funkciou \ac{mds}, ktorá vyžaduje ich prepojenie synchronizačnými káblami. S počítačom komunikujú cez sieťové rozhranie Ethernet (1GbE) prostredníctvom dátového servera softvéru LabOne. V jazyku Python ich ovládame pomocou knižnice zhinst-toolkit\footfullcite{zhinst_toolkit}, ktorá umožňuje nastavovať parametre prístroja (napr. časovú konštantu a rád filtra demodulátora či rozsah prúdového vstupu) a čítať namerané dáta.

\newpage
\section{Voltampérová charakteristika}
\subsection{Teoretický úvod}
\subsection{Použité meracie prístroje}
\subsection{Automatizácia merania a používateľské rozhranie}
\subsection{Praktická ukážka}
\subsection{Sumár merania}

\newpage
\section{Responzivita fotodetektorov}
\subsection{Teoretický úvod}
\subsection{Použité meracie prístroje}
\subsection{Automatizácia merania a používateľské rozhranie}
\subsection{Praktická ukážka}
\subsection{Sumár merania}

\newpage
\section{CPM -- Metóda konštantného fotoprúdu}
\subsection{Teoretický úvod}
\subsection{Použité meracie prístroje}
\subsection{Automatizácia merania a používateľské rozhranie}
\subsection{Praktická ukážka}
\subsection{Sumár merania}


% -- ZÁVER --
\newpage
\section{Záver}
% Sem príde zhrnutie, zovšeobecnenie, hodnotenie, prínos, návrh riešení a odporúčania do budúcnosti.

% -- BIBLIOGRAFIA --
\newpage
\printbibliography[title={Zoznam použitej literatúry}]

% -- PRÍLOHY --
\newpage
\section*{Prílohy}
\addcontentsline{toc}{section}{Prílohy} 

\subsection*{Príloha A: Schéma zapojenia prístrojov}
\begin{figure}[htbp]
    \centering
    % Odkomentuj tento riadok a uprav názov súboru, keď nahráš obrázok
    % \includegraphics[width=0.8\textwidth]{schema_zapojenia.png}
    \caption{Detailná schéma prepojenia pikoampérmetra a Lock-In zosilňovača.}
    \label{fig:schema_zapojenia}
\end{figure}

\end{document}